\subsubsection*{推论 2.2.}
\underline{在上述假设下, 若 $n(X/S,x)$ 为 $X/S$ 在 $x$ 处的消灭圈总数（1.8）, 则}
\[
\sum_{x\in F}\mu(X/S,x)=\sum_{x\in F}n(X/S,x).
\]

$f$ 的 Leray 谱序列给出
\begin{align*}
\chi(X)
&=\sum_{p,q}(-1)^{p+q}
\dim_{\mathbb Z/\ell}H^p(S,R^qf_*\mathbb Z/\ell)\\
&=\sum_q(-1)^q\chi(S,R^qf_*\mathbb Z/\ell).
\end{align*}
对 $S$ 上的每个可构造 $\mathbb Z/\ell$-模层 $G$, 有（[4] 或 SGA 5, X, 7.2）
\[
\chi(S,G)=\chi(S)\dim_{\mathbb Z/\ell}(G_{\bar\eta})
-\sum_{s\in S}\varepsilon_s(G),
\]
其中, 对每个 $s\in S$, 将 $\bar\eta$ 与 $S$ 在 $s$ 处的 Hensel 化的分式域的代数闭包等同后, 有\footnote{英文底本编校注: 原印本在 Swan 项前印为负号. 前一句引用的标准 Artin 导子公式, 以及原印本中将交错和等同于\emph{总}维数 $\dim+\dim\operatorname{Sw}$ 的等式 (c), 均要求此处为正号.}
\[
\varepsilon_s(G)=
\dim(G_{\bar\eta})-\dim(G_s)+\dim\operatorname{Sw}(G_{\bar\eta}).
\]

对 $s\in S$, 正合列 XIII, 2.4.6.3,
\[
\cdots\longrightarrow(R^qf_*\mathbb Z/\ell)_s
\longrightarrow(R^qf_*\mathbb Z/\ell)_{\bar\eta}
\longrightarrow\bigoplus_{x\in F(s)}
\Phi^q_{\bar\eta}(\mathbb Z/\ell)_x
\longrightarrow\cdots
\]
给出
\begin{align*}
&\sum_q(-1)^q\left(
\dim(R^qf_*\mathbb Z/\ell)_{\bar\eta}
-\dim(R^qf_*\mathbb Z/\ell)_s\right)\\
&\hspace{4em}=\sum_{x\in F(s)}\sum_q(-1)^q
\dim\Phi^q_{\bar\eta}(\mathbb Z/\ell)_x,
\end{align*}
及
\[
\sum_q(-1)^q\dim\operatorname{Sw}
(R^qf_*\mathbb Z/\ell)_{\bar\eta}
=\sum_{x\in F(s)}\sum_q(-1)^q
\dim\operatorname{Sw}\Phi^q_{\bar\eta}(\mathbb Z/\ell)_x.
\]
所以
\begin{equation}\tag{c}
\sum_q(-1)^q\varepsilon_s(R^qf_*\mathbb Z/\ell)
=\sum_{x\in F(s)}\dim\operatorname{tot}\Phi_x.
\end{equation}
于是 Euler--Poincar\'e 公式成为
\begin{align*}
\chi(X)
&=\sum_q(-1)^q\chi(S,R^qf_*\mathbb Z/\ell)\\
&=\sum_q(-1)^q\left(
\chi(S)\dim(R^qf_*\mathbb Z/\ell)_{\bar\eta}
-\sum_s\varepsilon_s(R^qf_*\mathbb Z/\ell)\right)\\
&=\chi(S)\sum_q(-1)^q
\dim(R^qf_*\mathbb Z/\ell)_{\bar\eta}
-\sum_s\sum_q(-1)^q\varepsilon_s(R^qf_*\mathbb Z/\ell)\\
&=\chi(S)\chi(X_{\bar\eta})
-\sum_s\sum_{x\in F(s)}\dim\operatorname{tot}\Phi_x\\
&=\chi(S)\chi(X_{\bar\eta})
-(-1)^n\sum_{x\in F}n(X/S,x).
\end{align*}
与 (2.1.1) 比较即证明 2.2.
